\code{FollowFork} is used to configure ProcControlAPI's behavior when the
associated Process forks.

\begin{apient}
enum follow_t {None, ImmediateDetach, DisableBreakpointsDetach, Follow }
\end{apient}

\apidesc{
The follow\_t type represents the configurable behaviors under forking.
None is specified when fork tracking is not available for the current platform.
ImmediateDetach means that forked children are never attached to. 
DisableBreakpointsDetach means that inherited breakpoints are removed from
forked children, and then the children are detached.
Follow is the default behavior, and it means that forked children are attached
to and remain under full control of ProcControlAPI.
}

\begin{apient}
static void setDefaultFollowFork(follow_t f)
\end{apient}

\apidesc{
This function sets the default forking behavior across all Process objects to f.
}

\begin{apient}
static follow_t getDefaultFollowFork()
\end{apient}

\apidesc{
This function returns the current default forking behavior across all Process
objects.
}

\begin{apient}
bool setFollowFork(follow_t f) const
\end{apient}

\apidesc{
This function sets the forking behavior for the associated Process object to f.
This function returns true on success and false on failure.
}

\begin{apient}
follow_t getFollowFork() const
\end{apient}

\apidesc{
This function returns the current forking behavior for the associated Process
object.
}